Les boies marines s'utilitzen sovint per a mesurar l'alçària de l'onatge. Una d'aquestes boies es mou seguint una oscil·lació harmònica de 3,00~m d'amplitud i 0,10~Hz de freqüència i l'ona es propaga a una velocitat de $0,50\ \mathrm{m\,s^{-1}}$.

\begin{apartat}{1}
Calculeu la longitud d'ona i el nombre d'ona.
\end{apartat}

\begin{apartat}{1}
Escriviu l'equació de les ones que fan moure la boia suposant que la fase inicial és zero.
\end{apartat}
